% !TEX program = xelatex
\RequirePackage{fix-cm}
\documentclass[12pt,oneside]{book}

% ==============================================================================
% TEMPLATE BÁO CÁO KHOA HỌC / LUẬN VĂN ĐẠI HỌC CẦN THƠ (CHUẨN QĐ 1799 - XeLaTeX)
% ==============================================================================

\usepackage{fancyhdr}
\usepackage{lipsum}
\usepackage{array}
\usepackage{indentfirst}
\usepackage{appendix}
\usepackage{subcaption}
\usepackage{longtable}
\usepackage[main=english,vietnamese]{babel}
\usepackage[a4paper,left=3.0cm,right=2.0cm,top=2.0cm,bottom=2.0cm,headheight=15pt,headsep=1.0cm,footskip=1.0cm]{geometry}
\usepackage{xcolor}
\usepackage{stmaryrd}
\usepackage{amsmath}
\usepackage{newtxmath}
\usepackage{fontspec}
\setmainfont{Times New Roman}
\setsansfont{Arial}
\setmonofont{Courier New}
\usepackage{graphicx}
\usepackage{setspace}
\usepackage{caption}

% Ngăn chặn hoàn toàn việc chèn trang trắng rỗng giữa các chương
\let\cleardoublepage\clearpage

\DeclareCaptionFont{size12}{\fontsize{12bp}{14bp}\selectfont}
% Tùy chỉnh định dạng của chú thích của Hình và Bảng
\captionsetup[table]{
  font={size12},
  labelfont=normalfont,
  textfont=normalfont,
  labelsep=colon,
  justification=raggedright,
  singlelinecheck=false,
  skip=6pt
}
\captionsetup[figure]{
  font={size12},
  labelfont=normalfont,
  textfont=normalfont,
  labelsep=colon,
  justification=centering,
  skip=6pt
}

\usepackage{float}
\usepackage{multirow}
\usepackage{minitoc}
\usepackage{tikz}
\usepackage{enumitem}
\setlist[enumerate]{itemsep=0pt,parsep=0pt,topsep=0pt,partopsep=0pt}
\setlist[itemize]{itemsep=0pt,parsep=0pt,topsep=0pt,partopsep=0pt}
\setlist[enumerate,1]{labelindent=1.0cm}
\setlist[itemize,1]{labelindent=1.0cm}
\usepackage{epigraph}
\usepackage{lscape}
\usepackage{titletoc}
\usepackage{booktabs}
\usepackage{arydshln}
\usepackage[nonumberlist]{glossaries}
\usepackage{calc}
\usepackage{titlesec}

% Cấu hình biblatex chuẩn numeric
\usepackage[
    backend=bibtex,
    style=numeric,
    backref=true,
    backrefstyle=three,
]{biblatex}

\addbibresource{references.bib}

\definecolor{linkColor}{HTML}{1B365D}
\usepackage[hidelinks,colorlinks=true,citecolor=linkColor,linkcolor=black,urlcolor=blue]{hyperref}

% Cấu hình cỡ chữ 13bp (tương đương 13pt Word với font Times New Roman từ fontspec)
\makeatletter
\DeclareRobustCommand\normalsize{%
   \@setfontsize\normalsize{13bp}{15bp}%
   \abovedisplayskip 12\p@ \@plus3\p@ \@minus7\p@
   \abovedisplayshortskip \z@ \@plus3\p@
   \belowdisplayshortskip 6.5\p@ \@plus3.5\p@ \@minus3\p@
   \belowdisplayskip \abovedisplayskip
   \let\@listi\@listI}
\let\@normalsize\normalsize
\makeatother
\AtBeginDocument{\normalsize}

% Thiết lập phân cấp đề mục
\setcounter{tocdepth}{2}
\setcounter{secnumdepth}{3}
\titleformat{\section}[block]{\fontsize{13bp}{15bp}\selectfont\bfseries}{\hspace{1.0cm}\thesection}{0.5em}{}
\titleformat{\subsection}[block]{\fontsize{13bp}{15bp}\selectfont\bfseries}{\hspace{1.0cm}\thesubsection}{0.5em}{}
\titleformat{\subsubsection}[block]{\fontsize{13bp}{15bp}\selectfont\bfseries}{\hspace{1.0cm}\thesubsubsection}{0.5em}{}

\usepackage[linesnumbered,ruled,vlined]{algorithm2e}
\usepackage{listings}

% ==============================================================================
% ĐỊNH DẠNG TIÊU ĐỀ CHƯƠNG VÀ MINITOC
% ==============================================================================

\definecolor{yourcolor}{HTML}{8a0e19}

\titleformat{\chapter}[display]
  {\normalfont\fontsize{14bp}{17bp}\selectfont\bfseries\centering}
  {\MakeUppercase{\chaptertitlename}\ \thechapter}
  {0.5ex}
  {\fontsize{14bp}{17bp}\selectfont\bfseries\MakeUppercase}
  []
\titlespacing*{\chapter}{0pt}{-20pt}{20pt}
\titleformat{name=\chapter,numberless}[display]
  {\normalfont\fontsize{14bp}{17bp}\selectfont\bfseries\centering}
  {}
  {0ex}
  {\fontsize{14bp}{17bp}\selectfont\bfseries\MakeUppercase}
  []

% Lệnh in đoạn tóm lược / mục lục nhỏ đầu chương
\newcommand{\printmyminitoc}[1]{%
	 #1\par\vspace*{1ex}
}

% ==============================================================================
% ĐỒNG BỘ KHOẢNG CÁCH VÀ TÊN DANH SÁCH BẢNG / HÌNH / BẢNG BIỂU THEO QĐ 1799
% ==============================================================================
\usepackage{etoolbox}
\AtBeginEnvironment{table}{\singlespacing}
\AtBeginEnvironment{longtable}{\singlespacing}
\makeatletter
% Vô hiệu hóa khoảng trắng 10pt chèn tự động giữa các chương trong .lof và .lot
\pretocmd{\@starttoc}{\begingroup\ifstrequal{#1}{lof}{\renewcommand{\addvspace}[1]{}}{}\ifstrequal{#1}{lot}{\renewcommand{\addvspace}[1]{}}{}}{}{}
\apptocmd{\@starttoc}{\endgroup}{}{}
\makeatother

% Đổi tiêu đề "Danh sách hình vẽ" thành "Danh sách hình" (Mục 1.3.1.9 QĐ 1799)
\addto\captionsvietnamese{%
  \renewcommand{\listfigurename}{Danh sách hình}%
  \renewcommand{\listtablename}{Danh sách bảng}%
}

\newcommand{\HRule}{\rule{\linewidth}{0.7mm}}
\newcommand{\Hrule}{\rule{\linewidth}{0.3mm}}

% ==============================================================================
% BẮT ĐẦU TÀI LIỆU CHÍNH
% ==============================================================================

\begin{document}
\selectlanguage{vietnamese}

% --- CỤM TIỀN TỐ 00_frontmatter: TRANG BÌA CHÍNH & BÌA PHỤ ---
% ==============================================================================
% TRANG BÌA CHÍNH (FRONTMATTER) - 1 TRANG DUY NHẤT, KHÔNG DƯ THỪA TRANG TRẮNG
% TIỀN TỐ: front_01_
% ==============================================================================

\begin{titlepage}
  \thispagestyle{empty}
  \begin{tikzpicture}[remember picture,overlay,inner sep=0,outer sep=0]
    \draw[black,line width=2.5pt]
      ([xshift=2.3cm,yshift=-1.4cm]current page.north west)
      rectangle
      ([xshift=-1.3cm,yshift=1.4cm]current page.south east);
  \end{tikzpicture}

  \begin{center}
    {\fontsize{14bp}{17bp}\selectfont \textbf{ĐẠI HỌC CẦN THƠ}\\
    \textbf{TRƯỜNG CÔNG NGHỆ THÔNG TIN VÀ TRUYỀN THÔNG}\par}
    \vspace{1.0cm}

    \includegraphics[height=3.0cm]{images/logoCTU.png}\\[1.0cm]

    {\fontsize{14bp}{17bp}\selectfont \textbf{BÁO CÁO KHOA HỌC / LUẬN VĂN THẠC SĨ}\\
    \textbf{HỌC PHẦN: [TÊN HỌC PHẦN / MÃ HỌC PHẦN]}\par}
    \vspace{2.6cm}

    {\fontsize{20bp}{25bp}\selectfont \bfseries [NHẬP TÊN ĐỀ TÀI TIẾNG VIỆT IN HOA TẠI ĐÂY TRONG PHẠM VI TỪ HAI ĐẾN BỐN DÒNG CÂN ĐỐI]\par}
    \vspace{0.4cm}

    {\fontsize{18bp}{22.5bp}\selectfont [ENTER THE ENGLISH TITLE IN UPPERCASE HERE MATCHING THE VIETNAMESE RESEARCH TOPIC]\par}
    \vspace{0.4cm}

    \vfill

    {\fontsize{14bp}{17bp}\selectfont
    \begin{tabular}{ll}
      \textbf{Giảng viên hướng dẫn:} & \textbf{[Học hàm, Học vị. Họ và Tên]} \\[0.35cm]
      \textbf{Học viên thực hiện:}   & \textbf{[Họ và Tên Học viên 1]} (MSHV: \textbf{[Mã số]}) \\
                                     & \textbf{[Họ và Tên Học viên 2]} (MSHV: \textbf{[Mã số]}) \\[0.2cm]
      \textbf{Chuyên ngành:}         & [Tên chuyên ngành] \\
      \textbf{Lớp:}                  & [Tên lớp / Khóa]
    \end{tabular}}

    \vfill

    {\fontsize{14bp}{17bp}\selectfont Cần Thơ, Tháng [MM]/[YYYY]}
  \end{center}
\end{titlepage}

% ==============================================================================
% TRANG BÌA PHỤ (FRONTMATTER) - SAU BÌA CHÍNH, KHÔNG LOGO
% TIỀN TỐ: front_01b_
% ==============================================================================

\begin{titlepage}
  \thispagestyle{empty}
  \begin{tikzpicture}[remember picture,overlay,inner sep=0,outer sep=0]
    \draw[black,line width=2.5pt]
      ([xshift=2.3cm,yshift=-1.4cm]current page.north west)
      rectangle
      ([xshift=-1.3cm,yshift=1.4cm]current page.south east);
  \end{tikzpicture}

  \begin{center}
    {\fontsize{14bp}{17bp}\selectfont \textbf{ĐẠI HỌC CẦN THƠ}\\
    \textbf{TRƯỜNG CÔNG NGHỆ THÔNG TIN VÀ TRUYỀN THÔNG}\par}
    \vspace{2.2cm}

    {\fontsize{14bp}{17bp}\selectfont \textbf{BÁO CÁO KHOA HỌC / LUẬN VĂN THẠC SĨ}\\
    \textbf{HỌC PHẦN: [TÊN HỌC PHẦN / MÃ HỌC PHẦN]}\par}

    \vfill

    {\fontsize{20bp}{25bp}\selectfont \bfseries [NHẬP TÊN ĐỀ TÀI TIẾNG VIỆT IN HOA TẠI ĐÂY TRONG PHẠM VI TỪ HAI ĐẾN BỐN DÒNG CÂN ĐỐI]\par}
    \vspace{0.4cm}

    {\fontsize{18bp}{22.5bp}\selectfont [ENTER THE ENGLISH TITLE IN UPPERCASE HERE MATCHING THE VIETNAMESE RESEARCH TOPIC]\par}
    \vspace{0.4cm}

    \vfill

    {\fontsize{14bp}{17bp}\selectfont
    \begin{tabular}{ll}
      \textbf{Giảng viên hướng dẫn:} & \textbf{[Học hàm, Học vị. Họ và Tên]} \\[0.35cm]
      \textbf{Học viên thực hiện:}   & \textbf{[Họ và Tên Học viên 1]} (MSHV: \textbf{[Mã số]}) \\
                                     & \textbf{[Họ và Tên Học viên 2]} (MSHV: \textbf{[Mã số]}) \\[0.2cm]
      \textbf{Chuyên ngành:}         & [Tên chuyên ngành] \\
      \textbf{Lớp:}                  & [Tên lớp / Khóa]
    \end{tabular}}

    \vfill

    {\fontsize{14bp}{17bp}\selectfont Cần Thơ, Tháng [MM]/[YYYY]}
  \end{center}
\end{titlepage}


% Cấu hình header & footer cho định dạng oneside theo QĐ 1799
\pagestyle{fancy}
\fancyhf{} % Xóa sạch toàn bộ header và footer mặc định
\renewcommand{\headrulewidth}{0pt} % QĐ 1799 không dùng đường kẻ trên đầu trang
\fancyfoot[C]{\fontsize{13bp}{15bp}\selectfont\thepage} % Số trang ở GIỮA PHÍA DƯỚI, cỡ 13

\fancypagestyle{plain}{
  \fancyhf{}
  \renewcommand{\headrulewidth}{0pt}
  \fancyfoot[C]{\fontsize{13bp}{15bp}\selectfont\thepage}
}

% --- CỤM TIỀN TỐ 00_frontmatter: THỦ TỤC TRƯỚC NỘI DUNG CHÍNH ---
\frontmatter
\normalsize
\setstretch{1.2}                     % Dãn dòng chuẩn 1.2 theo Mục 1.2.3 QĐ 1799
\setlength{\parskip}{6bp}            % Khoảng cách giữa các đoạn 6bp (before 6pt, after 0pt)
\setlength{\parindent}{1.0cm}        % Thụt đầu dòng đoạn văn 1.0cm
\titlespacing*{\section}{0pt}{0pt}{0pt}
\titlespacing*{\subsection}{0pt}{0pt}{0pt}
\titlespacing*{\subsubsection}{0pt}{0pt}{0pt}

% Bỏ dấu % ở các dòng dưới nếu cần dùng Trang xác nhận, Lời cảm ơn hoặc Lời cam đoan
% % ==============================================================================
% PHIẾU XÁC NHẬN VÀ PHÊ DUYỆT BÁO CÁO (FRONTMATTER)
% TIỀN TỐ: front_02_
% ==============================================================================

\noindent
\begin{minipage}[t]{0.48\textwidth}
\raggedright
{\fontsize{11bp}{14bp}\selectfont
ĐẠI HỌC CẦN THƠ\\[0.3ex]
\textbf{TRƯỜNG CNTT \& TRUYỀN THÔNG}\\[-1.2ex]
\hspace*{0.5cm}\rule{3.2cm}{0.4pt}
}
\end{minipage}
\hfill
\begin{minipage}[t]{0.50\textwidth}
\raggedleft
{\fontsize{11bp}{14bp}\selectfont
\begin{center}
\textbf{CỘNG HÒA XÃ HỘI CHỦ NGHĨA VIỆT NAM\\[0.3ex]
Độc lập -- Tự do -- Hạnh phúc}\\
\noindent\rule{4cm}{0.4pt}
\end{center}
}
\end{minipage}

\vspace{1.5cm}

\begin{center}
  {\fontsize{14bp}{17bp}\selectfont \textbf{PHIẾU ĐÁNH GIÁ VÀ XÁC NHẬN KẾT QUẢ}}\\
  {\fontsize{13bp}{16bp}\selectfont \textbf{BÁO CÁO KHOA HỌC / HỌC PHẦN [TÊN HỌC PHẦN]}}
\end{center}

\vspace{0.8cm}

\noindent\textbf{Tên đề tài:} [Nhập tên đề tài tiếng Việt tại đây].

\vspace{0.4cm}
\noindent\textbf{Danh sách học viên thực hiện:}
\begin{enumerate}[leftmargin=2.5cm,itemsep=1pt]
  \item \textbf{[Họ và Tên 1]} \hspace{1cm} MSHV: \textbf{[Mã số]}
  \item \textbf{[Họ và Tên 2]} \hspace{1cm} MSHV: \textbf{[Mã số]}
\end{enumerate}
\noindent\textbf{Chuyên ngành:} [Tên chuyên ngành] \hspace{2.5cm} \textbf{Lớp:} [Tên lớp]\\
\textbf{Thời gian báo cáo:} Tháng [MM] năm [YYYY]

\vspace{0.6cm}
\noindent\textbf{KẾT QUẢ ĐÁNH GIÁ CỦA HỘI ĐỒNG / GIẢNG VIÊN HƯỚNG DẪN:}
\begin{enumerate}[leftmargin=*,itemsep=2pt]
  \item Tiêu chí đánh giá 1: ................................................................. / ....... điểm
  \item Tiêu chí đánh giá 2: ................................................................. / ....... điểm
  \item Tiêu chí đánh giá 3: ................................................................. / ....... điểm
\end{enumerate}

\vspace{0.3cm}
\noindent\textbf{Tổng điểm đánh giá:} .............../ 100 điểm \hspace{2cm} \textbf{Xếp loại:} ....................................

\vspace{0.5cm}
\noindent\textbf{Ý kiến nhận xét:}\\
\noindent\dotfill\\
\noindent\dotfill\\
\noindent\dotfill

\vspace{1.2cm}

\noindent
\begin{minipage}[t]{0.45\textwidth}
  \centering
  \textbf{Đại diện học viên}\\
  \textit{(Ký và ghi rõ họ tên)}\\[2cm]
  \textbf{[Họ và Tên]}
\end{minipage}
\hfill
\begin{minipage}[t]{0.50\textwidth}
  \centering
  \textit{Cần Thơ, ngày ..... tháng ..... năm [YYYY]}\\
  \textbf{Giảng viên hướng dẫn / Chủ tịch Hội đồng}\\
  \textit{(Ký và ghi rõ họ tên)}\\[2cm]
  \textbf{[Học hàm, Học vị. Họ và Tên]}
\end{minipage}
          % 1. Trang xác nhận HĐ
% % ==============================================================================
% LỜI CẢM ƠN (FRONTMATTER)
% TIỀN TỐ: front_03_
% ==============================================================================

\chapter*{Lời cảm ơn}

Xin bày tỏ lòng biết ơn sâu sắc đến Quý Thầy Cô giảng dạy và hướng dẫn thuộc Trường Công nghệ Thông tin và Truyền thông, Đại học Cần Thơ đã tận tình chỉ bảo, đóng góp nhiều ý kiến chuyên môn quý báu giúp đề tài được hoàn thiện đúng tiến độ.

Xin gửi lời cảm ơn chân thành đến Ban Giám hiệu Nhà trường, Ban Chủ nhiệm Khoa cùng các bạn học viên và đồng nghiệp đã tạo mọi điều kiện thuận lợi về học liệu, cơ sở vật chất và môi trường nghiên cứu trong suốt thời gian thực hiện.

\vspace{1cm}
\begin{flushright}
  \textit{Cần Thơ, tháng [MM] năm [YYYY]}\\[0.5em]
  \textbf{[Tác giả / Nhóm thực hiện]}
\end{flushright}
            % 2. Lời cảm ơn
% ==============================================================================
% TÓM TẮT TIẾNG VIỆT (FRONTMATTER)
% TIỀN TỐ: front_04_
% ==============================================================================

\chapter*{Tóm tắt}

[Đoạn 1 -- Bối cảnh và Mục tiêu nghiên cứu: Trình bày ngắn gọn bối cảnh bài toán, tính cấp thiết và mục tiêu tổng quát của báo cáo hoặc luận văn (khoảng 200--350 từ theo Mục 1.3.1.3 QĐ 1799).]

[Đoạn 2 -- Phương pháp, Kết quả chính và Kết luận: Tóm lược phương pháp nghiên cứu, kỹ thuật thực nghiệm đã triển khai, các kết quả nổi bật đạt được và ý nghĩa thực tiễn hoặc hàm ý khoa học của công trình.]

\vspace{0.8cm}
\noindent\textbf{Từ khóa:} \textit{[Từ khóa 1, Từ khóa 2, Từ khóa 3, Từ khóa 4, Từ khóa 5]}
   % 3. Tóm tắt tiếng Việt
% ==============================================================================
% ABSTRACT TIẾNG ANH (FRONTMATTER)
% TIỀN TỐ: front_05_
% ==============================================================================

\chapter*{Abstract}

[Paragraph 1 -- Background and Objectives: Briefly introduce the research context, motivation, and primary objectives of the study (typically 200--350 words across 1--2 paragraphs).]

[Paragraph 2 -- Methodology, Key Findings, and Conclusions: Summarize the proposed methodology, experimental setup, main quantitative or qualitative findings, and practical implications.]

\vspace{0.8cm}
\noindent\textbf{Keywords:} \textit{[Keyword 1, Keyword 2, Keyword 3, Keyword 4, Keyword 5]}
   % 4. Tóm tắt tiếng Anh
% % ==============================================================================
% LỜI CAM ĐOAN KẾT QUẢ NGHIÊN CỨU (FRONTMATTER)
% TIỀN TỐ: front_06_
% ==============================================================================

\chapter*{Lời cam đoan}

Tôi/Chúng tôi xin cam đoan rằng công trình nghiên cứu này là kết quả thực hiện trung thực của bản thân/nhóm tác giả dưới sự hướng dẫn khoa học của [Học hàm, Học vị. Họ và Tên Giảng viên hướng dẫn]:

\begin{enumerate}[leftmargin=*]
  \item Toàn bộ nội dung, số liệu thu thập, mã nguồn và kết quả thực nghiệm trong tài liệu đều khách quan, trung thực và chưa từng được công bố hoặc sử dụng để xét duyệt ở bất kỳ luận văn, báo cáo cùng cấp nào khác.
  \item Tất cả tài liệu tham khảo, lý thuyết, biểu đồ hoặc đoạn mã trích dẫn từ các nguồn khác đều đã được viện dẫn rõ ràng và liệt kê đầy đủ trong Danh mục Tài liệu tham khảo.
\end{enumerate}

\vspace{1cm}
\begin{flushright}
  \textit{Cần Thơ, ngày ..... tháng ..... năm [YYYY]}\\[0.5em]
  \textbf{Người cam đoan}\\[1.5cm]
  \textbf{[Họ và Tên]}
\end{flushright}
          % 5. Lời cam đoan kết quả

\clearpage
\phantomsection
\addcontentsline{toc}{chapter}{Mục lục}
\tableofcontents

\clearpage
\phantomsection
\addcontentsline{toc}{chapter}{Danh sách Bảng}         % 7. Danh sách bảng (ĐẶT TRƯỚC HÌNH)
\listoftables

\clearpage
\phantomsection
\addcontentsline{toc}{chapter}{Danh sách Hình}         % 8. Danh sách hình (ĐẶT SAU BẢNG)
\listoffigures

% ==============================================================================
% DANH MỤC TỪ VIẾT TẮT (FRONTMATTER)
% TIỀN TỐ: front_09_
% ==============================================================================

\chapter*{Danh mục từ viết tắt}
\addcontentsline{toc}{chapter}{Danh mục từ viết tắt}

{\renewcommand{\baselinestretch}{1.0}\fontsize{12bp}{14bp}\selectfont
\begin{longtable}{lp{0.32\textwidth}p{0.44\textwidth}}
\toprule
\textbf{Viết tắt} & \textbf{Cụm từ tiếng Anh} & \textbf{Diễn giải tiếng Việt} \\ \midrule
\endfirsthead

\multicolumn{3}{l}{\textit{Tiếp tục từ trang trước}} \\ \toprule
\textbf{Viết tắt} & \textbf{Cụm từ tiếng Anh} & \textbf{Diễn giải tiếng Việt} \\ \midrule
\endhead

\midrule
\multicolumn{3}{r}{\textit{Xem trang kế tiếp}} \\
\endfoot

\bottomrule
\endlastfoot

AI  & Artificial Intelligence       & Trí tuệ nhân tạo \\
CTU & Can Tho University            & Trường Đại học Cần Thơ \\
LMS & Learning Management System    & Hệ thống quản trị học tập \\

\end{longtable}
}
    % 9. Danh mục từ viết tắt

% --- PHẦN NỘI DUNG CHÍNH (MAINMATTER) ---
\mainmatter
\setstretch{1.2}

% ==============================================================================
% CHƯƠNG 1: GIỚI THIỆU
% ==============================================================================

\chapter{Giới thiệu}\label{ch:gioi_thieu}

\startcontents[chapters]
\printmyminitoc{
Chương này giới thiệu bối cảnh nghiên cứu, tính cấp thiết của đề tài, câu hỏi nghiên cứu, mục tiêu và cấu trúc tổng thể của tài liệu.
}

% ==============================================================================
% MỤC 1.1 - 1.2: ĐẶT VẤN ĐỀ VÀ MỤC TIÊU NGHIÊN CỨU
% ==============================================================================

\section{Đặt vấn đề và Bối cảnh nghiên cứu}\label{sec:ch01_dat_van_de}

Phần văn bản chính sử dụng font Times New Roman cỡ chữ 13bp, giãn dòng 1.2 theo đúng Quyết định 1799 của Trường Đại học Cần Thơ \cite{ctu_qd1799}. Dòng đầu tiên của mỗi đoạn văn tự động thụt vào 1.0 cm và các đoạn cách nhau 6bp.

\subsection{Mục tiêu nghiên cứu}

Theo quy định trình bày, đề tài cần làm rõ mục tiêu tổng quát và các mục tiêu cụ thể \cite{sample_article2024}:

\subsubsection{Mục tiêu tổng quát}

Trình bày mục tiêu bao quát mà đề tài hoặc báo cáo khoa học hướng đến giải quyết trong phạm vi nghiên cứu đã xác định.

\subsubsection{Mục tiêu cụ thể}

Các gạch đầu dòng hoặc đánh số liệt kê sử dụng tùy chọn \texttt{[leftmargin=*]} để thẳng hàng đầu dòng 1.0 cm với các tiểu mục và không chèn thêm khoảng trống giữa các dòng liệt kê:
\begin{enumerate}[leftmargin=*]
  \item \textbf{Mục tiêu 1:} Hệ thống hóa cơ sở lý thuyết và các công trình nghiên cứu liên quan đến bài toán.
  \item \textbf{Mục tiêu 2:} Đề xuất quy trình, mô hình hoặc phương pháp giải quyết vấn đề nghiên cứu.
  \item \textbf{Mục tiêu 3:} Triển khai thực nghiệm, đánh giá kết quả định lượng và rút ra các nhận định khoa học.
\end{enumerate}

\section{Phạm vi và Đóng góp của đề tài}\label{sec:ch01_pham_vi}

Các nội dung liệt kê dạng gạch đầu dòng dùng môi trường \texttt{itemize}:
\begin{itemize}[leftmargin=*]
  \item \textbf{Phạm vi dữ liệu và thời gian:} Mô tả giới hạn tập dữ liệu, đối tượng khảo sát hoặc phạm vi thử nghiệm.
  \item \textbf{Đóng góp dự kiến:} Nêu rõ giá trị khoa học và khả năng ứng dụng thực tiễn của kết quả đạt được.
\end{itemize}

          % Chương 1: Giới thiệu
% ==============================================================================
% CHƯƠNG 2: TỔNG QUAN TÀI LIỆU
% ==============================================================================

\chapter{Tổng quan tài liệu}\label{ch:tong_quan}

\startcontents[chapters]
\printmyminitoc{
Chương này trình bày cơ sở lý thuyết nền tảng, mô hình toán học liên quan và tổng hợp các công trình nghiên cứu trước đây \cite{sample_book2022}.
}

\section{Cơ sở lý thuyết và Công thức toán học}\label{sec:ch02_co_so}

Các công thức toán học trong bài viết được đánh số tự động theo chương bằng môi trường \texttt{equation}. Ví dụ hàm mục tiêu chuẩn hóa:
\begin{equation}\label{eq:sample_loss}
  \mathcal{L}(\theta) = \frac{1}{N} \sum_{i=1}^{N} \left( y_i - f(x_i; \theta) \right)^2
\end{equation}
trong đó $N$ là số lượng mẫu dữ liệu, $y_i$ là giá trị thực tế và $f(x_i; \theta)$ là giá trị dự báo từ mô hình tại tham số $\theta$.

\section{Tổng hợp các công trình liên quan}\label{sec:ch02_lien_quan}

So sánh, phân tích ưu điểm và hạn chế của các nghiên cứu trước đây \cite{sample_article2024, sample_book2022} nhằm chỉ ra khoảng trống nghiên cứu mà đề tài sẽ tập trung giải quyết.
          % Chương 2: Tổng quan tài liệu
% ==============================================================================
% CHƯƠNG 3: PHƯƠNG PHÁP NGHIÊN CỨU
% ==============================================================================

\chapter{Phương pháp nghiên cứu}\label{ch:phuong_phap}

\startcontents[chapters]
\printmyminitoc{
Chương này mô tả quy trình nghiên cứu, thiết kế thực nghiệm cùng cách trình bày hình vẽ và bảng số liệu theo chuẩn QĐ 1799.
}

\section{Quy trình triển khai nghiên cứu}\label{sec:ch03_quy_trinh}

Hình vẽ và sơ đồ phải được đặt ngay sau đoạn mô tả (không đặt ngay sau tiêu đề mục), căn giữa và có chú thích cỡ 12bp ở phía dưới (ví dụ Hình \ref{fig:sample_pipeline}).

\begin{figure}[ht!]
  \centering
  \begin{tikzpicture}[>=latex]
    \node[draw, rectangle, rounded corners=3pt, fill=blue!8, minimum width=3.8cm, minimum height=0.9cm] (b1) at (0,0) {1. Thu thập dữ liệu};
    \node[draw, rectangle, rounded corners=3pt, fill=green!10, minimum width=3.8cm, minimum height=0.9cm] (b2) at (4.8,0) {2. Xử lý và Mô hình hóa};
    \node[draw, rectangle, rounded corners=3pt, fill=orange!12, minimum width=3.8cm, minimum height=0.9cm] (b3) at (9.6,0) {3. Đánh giá kết quả};
    \draw[->] (b1) -- (b2);
    \draw[->] (b2) -- (b3);
  \end{tikzpicture}
  \caption{Sơ đồ minh họa quy trình nghiên cứu tổng quát.}\label{fig:sample_pipeline}
\end{figure}

\section{Trình bày Bảng biểu theo quy chuẩn}\label{sec:ch03_bang}

Theo Phụ lục 7 của QĐ 1799, bảng không kẻ đường dọc, chỉ kẻ ngang ở đầu bảng, dưới tiêu đề cột và cuối bảng bằng gói \texttt{booktabs}; tiêu đề bảng đặt phía trên canh trái, cỡ chữ 12bp và nội dung bảng tự động giãn dòng 1.0 (như Bảng \ref{tab:sample_data}).

\begin{table}[ht!]
  \centering
  \caption{Bảng minh họa đặc tả tham số thực nghiệm chuẩn QĐ 1799.}\label{tab:sample_data}
  \begin{tabular}{llr}
    \toprule
    \textbf{Thành phần} & \textbf{Mô tả cấu hình} & \textbf{Giá trị} \\ \midrule
    Tập huấn luyện      & Tỷ lệ phân chia dữ liệu ban đầu & 70,0\% \\
    Tập kiểm thử        & Đánh giá độc lập cuối quy trình & 30,0\% \\
    Số vòng lặp tối đa  & Ngưỡng dừng hội tụ thuật toán & 100 \\
    \bottomrule
  \end{tabular}
\end{table}
        % Chương 3: Phương pháp nghiên cứu
% ==============================================================================
% CHƯƠNG 4: KẾT QUẢ VÀ THẢO LUẬN
% ==============================================================================

\chapter{Kết quả và Thảo luận}\label{ch:ket_qua}

\startcontents[chapters]
\printmyminitoc{
Chương này trình bày chi tiết kết quả thực nghiệm, phân tích số liệu và đối chiếu với các công trình trước đây.
}

\section{Kết quả thực nghiệm}\label{sec:ch04_ket_qua}

Trình bày kết quả đạt được kèm các bảng tổng hợp, biểu đồ minh họa và đánh giá thống kê cụ thể.

\section{Thảo luận}\label{sec:ch04_thao_luan}

Phân tích ý nghĩa khoa học của kết quả thu được, chỉ ra các nguyên nhân và thảo luận giới hạn của nghiên cứu.
            % Chương 4: Kết quả và thảo luận
% ==============================================================================
% CHƯƠNG 5: KẾT LUẬN VÀ ĐỀ XUẤT
% ==============================================================================

\chapter{Kết luận và Đề xuất}\label{ch:ket_luan}

\startcontents[chapters]
\printmyminitoc{
Chương này tổng kết các đóng góp chính của nghiên cứu và đề xuất hướng phát triển tiếp theo.
}

\section{Kết luận}\label{sec:ch05_ket_luan}

Tóm tắt các kết quả cốt lõi đạt được gắn với mục tiêu và câu hỏi nghiên cứu ban đầu.

\section{Đề xuất và Hướng phát triển}\label{sec:ch05_de_xuat}

Đưa ra các kiến nghị thực tiễn và hướng mở rộng nghiên cứu trong tương lai.
           % Chương 5: Kết luận và đề xuất

% --- TÀI LIỆU THAM KHẢO ---
\clearpage
\phantomsection
\addcontentsline{toc}{chapter}{Tài liệu tham khảo}
\begin{singlespace}
\printbibliography
\end{singlespace}

% --- CỤM TIỀN TỐ app_: PHỤ LỤC ---
% ==============================================================================
% PHỤ LỤC (APPENDICES)
% FILE CHỦ PHỤ LỤC: app_main.tex
% ==============================================================================

\clearpage
\begin{appendices}
  \singlespacing                                 % Dãn dòng Phụ lục = 1.0 theo Mục 1.2.3 QĐ 1799
  \renewcommand{\thechapter}{\arabic{chapter}}   % Đánh số bằng chữ số (1, 2, 3) thay vì (A, B, C)
  \renewcommand{\appendixname}{Phụ lục}
  \pagenumbering{arabic}                         % Đánh số trang Phụ lục lại từ 1
  \setcounter{page}{1}
  % Tiêu đề Phụ lục chữ in thường, in đậm, cỡ 14bp, canh giữa theo Mục 1.3.3.2
  \titleformat{\chapter}[display]
    {\normalfont\fontsize{14bp}{17bp}\selectfont\bfseries\centering}
    {Phụ lục\ \thechapter}
    {0.5ex}
    {\fontsize{14bp}{17bp}\selectfont\bfseries}
    []
  % ==============================================================================
% PHỤ LỤC 1: BIỂU MẪU KHẢO SÁT / TÀI LIỆU BỔ TRỢ
% ==============================================================================

\chapter{Biểu mẫu Khảo sát hoặc Bảng số liệu Bổ trợ}\label{app:bieu_mau}

Toàn bộ phần Phụ lục được trình bày ở mức giãn dòng 1.0 và đánh số trang lại từ 1 theo Mục 1.3.3.2 QĐ 1799.

\section*{Phần 1: Cấu trúc câu hỏi mẫu}
\begin{enumerate}[leftmargin=*]
  \item Câu hỏi khảo sát thứ nhất (giữ nguyên dạng nguyên bản nếu dùng để thu thập dữ liệu).
  \item Câu hỏi khảo sát thứ hai kèm các phương án lựa chọn:
  \begin{itemize}
    \item (A) Phương án 1
    \item (B) Phương án 2
    \item (C) Phương án 3
  \end{itemize}
\end{enumerate}

  % ==============================================================================
% PHỤ LỤC 2: MÃ NGUỒN / MINH CHỨNG TRIỂN KHAI
% ==============================================================================

\chapter{Minh chứng Triển khai và Tài nguyên Công khai}\label{app:minh_chung}

Liệt kê các đường dẫn kho mã nguồn, bộ dữ liệu hoặc kịch bản thực nghiệm bổ sung (nếu có):

\begin{enumerate}[leftmargin=*]
  \item \textbf{Kho mã nguồn (GitHub):}\\
  \url{https://github.com/username/project-repo}
  \item \textbf{Dữ liệu thực nghiệm:} Mô tả ngắn gọn vị trí lưu trữ và hướng dẫn tái lập kết quả.
\end{enumerate}

\end{appendices}


\end{document}
